% !TeX root = booklet-01.tex
\section{انتقال داده و خروجی سه‌حالته}
\label{sec:ch1-pass}
\subsection{کنترل مسیر با تولید سطح منطقی فرق دارد}
در گیت ایستای قبلی، ورودی به \lr{Gate}ها می‌رفت و خروجی از ریل تغذیه یا زمین
تولید می‌شد. در \lr{Pass Transistor}، خود داده از مسیر \lr{Source--Drain}
عبور می‌کند؛ \lr{Gate} فقط مسیر عبور را کنترل می‌کند.
این تفاوت توضیح می‌دهد چرا یک کلید روشن الزاماً ولتاژ داده را بدون تغییر عبور نمی‌دهد
\cite[بخش \secnum{۱.۴.۶}]{weste}.

\textbf{مثال آموزشی برای عبور یک.} \lr{nMOS} با \lr{Gate} در
\(V_{\mathrm{DD}}\) گره خروجی را از ولتاژ پایین شارژ می‌کند.
با بالا رفتن خروجی، پایانه کم‌ولتاژ مسیر نیز بالا می‌رود و
\(V_{gs}=V_{\mathrm{DD}}-V_{\mathrm{out}}\) کم می‌شود.
در مدل آستانه ثابت، حد شارژ وقتی است که این اختلاف به \(V_{tn}\) برسد:
\begin{equation}\label{eq:ch1-pass-high}
V_{\mathrm{out,lim}}=V_{\mathrm{DD}}-V_{tn}.
\end{equation}
از رابطه \(V_{\mathrm{DD}}-V_{\mathrm{out}}=V_{tn}\)،
ابتدا \(V_{\mathrm{out}}\) را به طرف دیگر و سپس \(V_{tn}\) را کم کردیم.
برای تغذیه فرضی \(\qty{1}{\volt}\) و آستانه \(\qty{0.3}{\volt}\)،
حد خروجی \(\qty{0.7}{\volt}\) است؛ نه یک ولت.
این حد در مدل بدون نشتی، بدون بار رسانشی و با شروع از ولتاژ پایین بیان شده است.
با \lr{Body Effect}، خود آستانه نیز تغییر می‌کند؛ تحلیل کامل به جزوه سوم تعلق دارد.

\lr{nMOS} در انتقال صفر به زمین چنین افت آستانه‌ای ایجاد نمی‌کند؛
ازاین‌رو کتاب از «صفر قوی و یک ضعیف» صحبت می‌کند.
\lr{pMOS} برعکس، یک را قوی و صفر را با افت انتقال می‌دهد:
با \lr{Gate} صفر، حد دشارژ در مدل مشابه نزدیک \(|V_{tp}|\) است
\cite[شکل \secnum{۱.۲۰}]{weste}.
«ضعیف» یعنی سطح از ریل فاصله دارد و توان راه‌اندازی نیز به شرایط مدار وابسته است؛
این اصطلاح به‌تنهایی تعیین نمی‌کند سطح برای ورودی گیت بعد معتبر است یا خیر.

\subsection{\lr{Transmission Gate}: ترکیب دو نوع مکمل}
با موازی کردن \lr{nMOS} و \lr{pMOS} و اعمال کنترل‌های مکمل،
\lr{Transmission Gate (TG)} می‌سازیم. برای \(E=1\)،
\lr{Gate} ترانزیستور \lr{nMOS} یک و \lr{Gate} ترانزیستور \lr{pMOS}
صفر است؛ مسیر داده فعال می‌شود. برای \(E=0\)، هر دو خاموش‌اند.
در انتقال صفر، \lr{nMOS} و در انتقال یک، \lr{pMOS} سطح مناسب را پشتیبانی می‌کند
\cite[شکل \secnum{۱.۲۱}]{weste}.
شکل \ref{fig:ch1-tg} این موازی بودن را نشان می‌دهد.

\begin{figure}[ht]
\centering
\begin{latin}
\begin{circuitikz}[american]
\draw (0,1) node[pmos,rotate=-90] (p) {};
\draw (0,-1) node[nmos,rotate=90] (n) {};
\draw (p.D)--(-2,1)--(-2,-1)--(n.D);
\draw (p.S)--(2,1)--(2,-1)--(n.S);
\draw (-2,0) to[short,*-o] (-3,0) node[left] {$a$};
\draw (2,0) to[short,*-o] (3,0) node[right] {$b$};
\draw (p.G)--(0,2.2) node[above] {$\overline{E}$};
\draw (n.G)--(0,-2.2) node[below] {$E$};
\node[right] at (0.7,1.5) {pMOS};\node[right] at (0.7,-1.5) {nMOS};
\end{circuitikz}
\end{latin}
\caption{بازترسیم آموزشی \lr{TG} بر پایه کتاب \cite{weste}.
پایانه‌های داده \(a,b\) دو سوی یک مسیر قابل‌کنترل‌اند؛
بدنه \lr{nMOS} به زمین و بدنه \lr{pMOS} به تغذیه متصل است.}
\label{fig:ch1-tg}
\end{figure}

مکمل کنترل رایگان نیست؛ اگر \(\overline{E}\) از قبل موجود نباشد،
یک \lr{Inverter} یا مدار تولید مکمل لازم است.
خود \lr{TG} دو ترانزیستور دارد، اما شمارش کل مدار باید هزینه کنترل را هم در نظر بگیرد.
برچسب \(a,b\) عمدی است؛ جهت انتقال می‌تواند با بایاس عوض شود.
این کاربرد برای ولتاژهای داخل ریل‌های مجاز بررسی می‌شود، نه یک کلید عمومی برای هر ولتاژی.

\subsection{تفاوت \lr{Restoring} و \lr{Nonrestoring}}
یک گیت \lr{Restoring}، سطح ورودی معتبر را با استفاده از تغذیه به خروجی نزدیک
ریل بازتولید می‌کند. مثلاً دو \lr{Inverter} پشت‌سرهم، داده را بدون تغییر منطقی
و با بازسازی سطح عبور می‌دهند؛ محدوده معتبر ورودی باید رعایت شود.
\lr{TG} داده را عبور می‌دهد و نویز یا افت ولتاژ آن را به‌خودی‌خود حذف نمی‌کند؛
بنابراین \lr{Nonrestoring} است \cite[بخش \secnum{۱.۴.۷}]{weste}.
«عبور هر دو سطح ریل» را نباید با «بازسازی ورودی خراب» برابر دانست.
زنجیره‌های بلند مسیر عبوری نیز مقاومت و ظرفیت دارند و در مدل واقعی تاخیرشان افزایش می‌یابد.

\subsection{\lr{Tristate}: مقدار \texorpdfstring{\(Z\)}{Z} یک سطح ولتاژ سوم نیست}
\lr{Tristate Buffer} در حالت فعال \(Y=A\) و در حالت غیرفعال خروجی
\lr{High Impedance} دارد. \(Z\) یعنی مدار خروجی را به‌طور فعال به صفر یا یک
نمی‌کشد؛ یک ولتاژ مشخص میان صفر و تغذیه نیست.
یک \lr{TG} همین جدول منطقی فعال/غیرفعال را دارد، ولی \lr{Restoring} نیست
\cite[صفحات چاپی ۱۴ و ۱۵]{weste}.

\lr{Tristate Inverter} در حالت فعال، خروجی \(\overline{A}\) می‌دهد و در حالت
غیرفعال \(Z\). برای ساخت آن، در مسیر بالایی \lr{Inverter} یک \lr{pMOS}
با کنترل \(\overline{E}\) و در مسیر پایینی یک \lr{nMOS} با کنترل \(E\)
سری می‌کنیم. پس:
\begin{equation}\label{eq:ch1-tristate}
F_{\mathrm{PUN}}=E\overline{A},\qquad
F_{\mathrm{PDN}}=EA.
\end{equation}
اگر \(E=1\)، شبکه‌ها مانند \lr{Inverter} عمل می‌کنند؛ اگر \(E=0\)،
هر دو قطع‌اند. جدول \ref{tab:ch1-tri} نتیجه را نشان می‌دهد.
شماتیک شکل \ref{fig:ch1-tri} محل ترانزیستورهای فعال‌ساز را مشخص می‌کند.

\begin{figure}[ht]
\centering
\begin{latin}
\begin{circuitikz}[american,scale=0.85,transform shape]
\draw (0,4.5) node[pmos] (pe) {} (0,2.7) node[pmos] (pa) {};
\draw (0,-0.3) node[nmos] (na) {} (0,-2.1) node[nmos] (ne) {};
\draw (pe.S)--(0,5.5) node[above] {$V_{\mathrm{DD}}$};
\draw (pe.D)--(pa.S) (pa.D)--(0,1.2)--(na.D) (na.S)--(ne.D);
\draw (ne.S)--(0,-3.1) node[ground] {};
\draw (0,1.2) to[short,*-o] (2,1.2) node[right] {$Y$};
\node[left] at (pe.G) {$\overline{E}$};\node[left] at (pa.G) {$A$};
\node[left] at (na.G) {$A$};\node[left] at (ne.G) {$E$};
\end{circuitikz}
\end{latin}
\caption{بازترسیم آموزشی \lr{Tristate Inverter} بر پایه شکل \secnum{۱.۲۷} کتاب
\cite{weste}. برچسب‌های یکسان \lr{Gate} یعنی ورودی مشترک؛
بدنه‌های \lr{pMOS} به تغذیه و \lr{nMOS} به زمین متصل‌اند.}
\label{fig:ch1-tri}
\end{figure}

\begin{table}[ht]
\centering
\caption{جدول \lr{Tristate Inverter}؛ مکمل کنترل در شمارش قطعات فرض نشده است.}
\label{tab:ch1-tri}
\begin{tabular}{ccc}
\toprule
\(E\)&\(A\)&\(Y\)\\
\midrule
۰&۰&\(Z\)\\
۰&۱&\(Z\)\\
۱&۰&۱\\
۱&۱&۰\\
\bottomrule
\end{tabular}
\end{table}
هسته این مدار چهار ترانزیستور دارد و خروجی فعالش \lr{Restoring} است.
برای \lr{Tristate Buffer} می‌توان یک \lr{Inverter} عادی پیش از آن گذاشت تا
دو معکوس‌سازی یکدیگر را خنثی کنند.

اگر چند فرستنده به یک \lr{Bus} وصل شوند، باید قرارداد فعال‌سازی مانع فرمان‌های
متضاد شود. دو فرستنده فعال با داده مخالف، \lr{Contention} ایجاد می‌کنند؛
نبود هیچ فرستنده فعال نیز می‌تواند \lr{Bus} را شناور کند.
حتی اگر دو داده در یک لحظه یکسان باشند، فعال‌سازی هم‌زمان درستی پروتکل را تضمین نمی‌کند.
تاخیر مسیرهای کنترل ممکن است هنگام جابه‌جایی فرستنده‌ها، هم‌پوشانی موقت بسازد؛
ازاین‌رو جایگزینی با \lr{Multiplexer} در بسیاری از طراحی‌ها مفید است
\cite[بخش \secnum{۱.۴.۷}]{weste}.

\section{\lr{Multiplexer}: انتخاب یکی از چند ورودی}
\label{sec:ch1-mux}
\subsection{از تعریف انتخاب تا رابطه منطقی}
\lr{Multiplexer} یا \lr{MUX} دوورودی، با انتخاب \(S=0\) داده \(D_0\)
و با \(S=1\) داده \(D_1\) را به خروجی می‌فرستد. بنابراین:
\begin{equation}\label{eq:ch1-mux}
Y=\overline{S}D_0+SD_1.
\end{equation}
وقتی \(S=0\)، جمله دوم صفر و جمله اول \(D_0\) است؛ وقتی \(S=1\)،
جمله اول صفر و جمله دوم \(D_1\). این جای‌گذاری ساده، هم رابطه و هم جهت انتخاب
را کنترل می‌کند \cite[بخش \secnum{۱.۴.۸}]{weste}.
زیرنویس‌های صفر و یک بخشی از نام دو ورودی‌اند؛ ضرب در صفر یا یک نیستند.

\subsection{پیاده‌سازی با دو \lr{TG}}
دو \lr{TG} را از \(D_0\) و \(D_1\) به گره مشترک \(Y\) وصل می‌کنیم.
کنترل فعال‌ساز مسیر اول \(\overline{S}\) و مسیر دوم \(S\) است.
برای هر مسیر، ترانزیستور \lr{nMOS} همان کنترل فعال‌ساز و \lr{pMOS} مکمل آن
را می‌گیرد. در مدل کنترل‌های مکمل ایدئال، دقیقاً یک مسیر باز است.
شکل \ref{fig:ch1-mux} نمای بلوکی همین اتصال را نشان می‌دهد.

\begin{figure}[ht]
\centering
\begin{latin}
\begin{tikzpicture}[x=1.2cm,y=1cm]
\node[draw,minimum width=1.8cm,minimum height=0.8cm] (t0) at (0,1.2) {TG};
\node[draw,minimum width=1.8cm,minimum height=0.8cm] (t1) at (0,-1.2) {TG};
\draw[->] (-2.2,1.2) node[left] {$D_0$}--(t0.west);
\draw[->] (-2.2,-1.2) node[left] {$D_1$}--(t1.west);
\draw (t0.east)--(2,1.2)--(2,-1.2)--(t1.east);
\draw[->] (2,0)--(3.1,0) node[right] {$Y$};\fill (2,0) circle (1.5pt);
\draw[->] (0,2.4) node[above] {$E_0=\overline{S}$}--(t0.north);
\draw[->] (0,-2.4) node[below] {$E_1=S$}--(t1.south);
\end{tikzpicture}
\end{latin}
\caption{نمای بلوکی آموزشی \lr{MUX} بر پایه شکل \secnum{۱.۲۸} کتاب \cite{weste}.
هر بلوک \lr{TG} شامل دو ترانزیستور شکل \ref{fig:ch1-tg} است؛
پیکان کنترل مسیر داده نیست.}
\label{fig:ch1-mux}
\end{figure}

هسته \lr{MUX} دو \lr{TG}، چهار ترانزیستور دارد؛ با تولید محلی
\(\overline{S}\) به کمک یک \lr{Inverter}، تعداد به شش می‌رسد.
این مدار \lr{Nonrestoring} است. قرار دادن \lr{Inverter} در خروجی،
سطح را بازسازی می‌کند ولی تابع را معکوس می‌کند؛ برای خروجی غیرمعکوس بازسازی‌شده،
به دو معکوس‌سازی یا پیاده‌سازی مناسب دیگری نیاز داریم.

در پیاده‌سازی گیتی، دو \lr{Tristate Inverter} با فعال‌سازهای مکمل
می‌توانند خروجی مشترک \(\overline{Y}\) را بسازند.
روش دیگر، گیت \lr{AOI22} با ورودی‌های \(\overline{S},D_0,S,D_1\) است.
کتاب این پیاده‌سازی‌ها را مقایسه می‌کند؛ استفاده از مکمل بودن دو کنترل،
گاهی اجازه ساده‌سازی شبکه را می‌دهد \cite[شکل \secnum{۱.۲۹}]{weste}.
این ساده‌سازی به قید کنترل‌ها متکی است و برای چهار ورودی مستقل قابل تعمیم نیست.

\subsection{انتخاب چهار داده و معنای دو بیت انتخاب}
یک \lr{MUX} چهارورودی را با سه \lr{MUX} دوورودی می‌سازیم.
در مرحله اول، یکی از \(D_0,D_1\) و یکی از \(D_2,D_3\) با \(S_0\)
انتخاب می‌شوند؛ در مرحله دوم، \(S_1\) یکی از دو نتیجه را انتخاب می‌کند.
به‌ترتیب برای \((S_1,S_0)=(0,0),(0,1),(1,0),(1,1)\)،
خروجی \(D_0,D_1,D_2,D_3\) است. شماره ورودی برابر \(2S_1+S_0\) می‌شود.
این ساختار دو مرحله مسیر داده دارد؛ نام چهارورودی به معنی چهار گیت سری نیست.

روش دیگر، چهار مسیر سه‌حالته با فعال‌سازهای
\(\overline{S_1}\overline{S_0}\)، \(\overline{S_1}S_0\)،
\(S_1\overline{S_0}\) و \(S_1S_0\) است
\cite[شکل \secnum{۱.۳۰}]{weste}.
برای ورودی پایدار دقیقاً یکی فعال است؛ کنترل گذرا همچنان باید مانع وصل هم‌زمان
مسیرهای داده متضاد شود. مدل منطقی به‌تنهایی زمان‌بندی واقعی کنترل‌ها را اثبات نمی‌کند.

\subsection{پرسش‌های انتقال و انتخاب}
\begin{enumerate}
\item \textbf{چرا یک \lr{nMOS} با کنترل یک، همیشه یک کامل عبور نمی‌دهد؟}
با بالا رفتن خروجی، اختلاف \(V_{gs}\) کم می‌شود و شرط هدایت از دست می‌رود.
\item \textbf{آیا \(Z\) همان صفر منطقی است؟}
خیر؛ به معنی نداشتن راه‌اندازی فعال است. یک خازن شناور می‌تواند برای مدتی یک قبلی را نگه دارد.
\item \textbf{آیا \lr{TG} نویز ورودی را حذف می‌کند؟}
خیر؛ عبوری و \lr{Nonrestoring} است. بازسازی به یک مدار دارای بهره و تغذیه مناسب نیاز دارد.
\item \textbf{در \lr{MUX} با \(S=1\)، تغییر \(D_0\) چه اثری بر خروجی منطقی دارد؟}
در مدل پایدار اثری ندارد؛ مسیر انتخاب‌شده \(D_1\) است. کوپل خازنی و گذرا در این حکم مدل نشده‌اند.
\end{enumerate}
